% ===================================================================
%  CUADRO N.º 9 — La Montaña y el Balneario. — El Bosque y la Caza.
%  Traduction 2026 d'apres texto/io, controlee sur texto/fr et
%  texto/en. Consigne : voir texto/es/10-cuadro-01.tex.
% ===================================================================

%%K t09-apar-1 apar
\VUsaut{4.0mm}
\VUtitre{15.0pt}{60}{CUADRO N.º 9}
\VUsaut{3.0mm}

%%K t09-tit-1 sub
\VUcentre{12.0pt}{0}{\textit{Primera escena.}}
\VUsaut{3.0mm}

%%K t09-tit-2 sub
\VUpk{11.4pt}{40}{La Montaña. --- La Ciudad balneario.}
\VUsaut{3.0mm}

%%K t09-c1-01-1 p
1. --- Llevo ocho días en Pratz. No puedo expresar toda la admiración que
sentí cuando me encontré ante la maravillosa \VUgras{cordillera}
\textsuperscript{(1)} de los Pirineos, cuya \VUgras{cresta} \textsuperscript{(2)} se
recorta sobre el fondo azul del cielo como un festón gigantesco.

%%K t09-c1-02-1 p
2. --- No imaginaba que las \VUgras{cumbres} \textsuperscript{(3)} pirenaicas
alcanzaran una \VUgras{altura} tan grande, y me quedé pasmado ante los
magníficos \VUgras{glaciares} \textsuperscript{(5)} y los \VUgras{picos}
\textsuperscript{(4)} cubiertos de nieve por los que ruedan \VUgras{aludes} tan
terribles.

%%K t09-tit-3 sub
\VUsaut{3.0mm}
\VUpk{10.2pt}{40}{Paisaje de montaña.}
\VUsaut{3.0mm}

%%K t09-c1-03-1 p
3. --- Al día siguiente mismo de mi llegada fui al viejo \VUgras{volcán}
\textsuperscript{(6)} apagado que hay cerca de Pratz. En los bordes áridos y
desnudos del \VUgras{cráter} se ven todavía muy bien las huellas del paso
de la \VUgras{lava} que corría, hace varios siglos, por la \VUgras{ladera de
la montaña} \textsuperscript{(7)}. Conseguí encontrar, en una de las
\VUgras{vertientes}, algunos fragmentos de esa lava.

%%K t09-c1-04-1 p
4. --- Pratz está en la frontera, y la \VUgras{fortaleza} \textsuperscript{(8)} que
defiende el \VUgras{puerto} \textsuperscript{(27)} que separa Francia de España
recuerda del todo a esos viejos castillos de las orillas del Rin que
parecen acechar una \VUgras{presa} desde lo alto de su roca.

%%K t09-c1-05-1 p
5. --- Un \VUgras{águila} \textsuperscript{(9)} enorme, que tiene su nido no lejos
del \VUgras{fuerte} \textsuperscript{(8)}, atrae a menudo mi atención. Esta
\VUgras{ave de rapiña}, de \VUgras{pico} poderoso, lleva con frecuencia a sus
crías un cordero o un cabrito sujeto en sus \VUgras{garras}. Sus
\VUgras{alas} son tan grandes que puede planear largo rato con su pesada
carga.

%%K t09-tit-4 sub
\VUsaut{3.0mm}
\VUpk{10.2pt}{40}{El Excursionista.}
\VUsaut{3.0mm}

%%K t09-c1-06-1 p
6. --- El paseo más agradable de aquí es la excursión a la \VUgras{cascada}
\textsuperscript{(10)} de Cau. El \VUgras{salto de agua} cae en remolinos y luego se
pierde en un bonito \VUgras{arroyo} en el \VUgras{pinar de abetos}
\textsuperscript{(11)} situado al oeste de la ciudad. Subimos por ese bosque hasta
el \VUgras{observatorio meteorológico} \textsuperscript{(12)} y, desde lo alto del
\VUgras{mirador}, abarcamos el \VUgras{panorama} entero de la comarca. El
\VUgras{paisaje} es soberbio y la \VUgras{naturaleza} parece haber prodigado
a este país todos los tesoros de que dispone.

%%K t09-c1-07-1 p
7. --- Para bajar tomamos el \VUgras{funicular} \textsuperscript{(13)} y nos paramos
en una pequeña \VUgras{estación} junto a las \VUgras{ruinas} \textsuperscript{(14)} de
un viejo \VUgras{castillo feudal} habitado antaño por los señores de Pratz.

%%K t09-c1-08-1 p
8. --- ¡Pobre castillo! ¿Qué pensará al ver a sus pies la coqueta
\VUgras{ciudad balneario} \textsuperscript{(15)}, que hoy es una
\VUgras{estación termal} de moda?

%%K t09-c1-08-2 p
El \VUgras{clima} es tan agradable y el \VUgras{agua mineral} tan benéfica
que la \VUgras{cura} que se sigue en el \VUgras{sanatorio} \textsuperscript{(17)} es
un verdadero placer.

%%K t09-tit-5 sub
\VUsaut{3.0mm}
\VUpk{10.2pt}{40}{Una agradable ciudad termal.}
\VUsaut{3.0mm}

%%K t09-c1-09-1 p
9. --- Por lo demás, se ha hecho todo para que la estancia sea agradable.
Se construyó un gran \VUgras{viaducto} \textsuperscript{(16)} para poder tender una
vía férrea; el \VUgras{parque} \textsuperscript{(18)} del \VUgras{balneario}, en el
que se dan \VUgras{conciertos}, está arreglado de un modo encantador. Se
han construido varios graciosos \VUgras{chalés} \textsuperscript{(19)} alrededor del
\VUgras{casino} \textsuperscript{(21)}, y una \VUgras{calzada} \textsuperscript{(22)} muy bien
cuidada hace facilísimas las comunicaciones.

%%K t09-c1-10-1 p
10. --- Un simple \VUgras{talud} \textsuperscript{(23)} separa la \VUgras{carretera}
\textsuperscript{(20)} del \VUgras{valle} \textsuperscript{(25)} propiamente dicho, en cuyo
fondo se extiende un bonito \VUgras{lago} \textsuperscript{(26)} que alimenta un
\VUgras{arroyo} \textsuperscript{(24)} muy claro.

%%K t09-c1-11-1 p
11. --- Al otro lado del valle se explota una \VUgras{mina} \textsuperscript{(28)} de
hierro. Fui varias veces a ver bajar a los \VUgras{mineros} por el
\VUgras{pozo}, y hasta entré en la \VUgras{galería} donde pasan el día
extrayendo el \VUgras{mineral} que contiene el \VUgras{metal}.

%%K t09-c1-12-1 p
12. --- Se construyó una importante \VUgras{fundición} junto a la mina para
aprovechar enseguida las riquezas que de ella se sacan. El \VUgras{alto
horno} \textsuperscript{(29)}, calentado con el \VUgras{carbón de piedra} que
abunda aquí, permite obtener \VUgras{hierro colado} de un modo rápido y
barato.

%%K t09-c1-13-1 p
13. --- No lejos de la mina se abre una \VUgras{cantera} \textsuperscript{(30)}. No es
ni \VUgras{granito}, ni \VUgras{pórfido}, ni siquiera \VUgras{arenisca} lo que
el viejo \VUgras{cantero} arranca con su \VUgras{pico}, sino simple
\VUgras{piedra de sillería} para construir casas.

%%K t09-c1-14-1 p
14. --- Uno de los más bellos adornos de este país es el \VUgras{abeto}
\textsuperscript{(31)}. En todas partes --- en el fondo de un \VUgras{barranco}
\textsuperscript{(32)}, a la orilla de un \VUgras{torrente} \textsuperscript{(33)}, junto a
una \VUgras{sima} \textsuperscript{(34)} o a un precipicio ---, sus finas agujas
ponen su nota verde.

%%K t09-tit-6 sub
\VUsaut{3.0mm}
\VUpk{10.2pt}{40}{Andar a pie.}
\VUsaut{3.0mm}

%%K t09-c1-15-1 p
15. --- Aprovecho mis vacaciones para dar largos paseos. Los
\VUgras{caminos} \textsuperscript{(35)} que serpentean por la montaña permiten
recorrer, sin darse cuenta de la distancia, varios \VUgras{kilómetros} y a
menudo varias \VUgras{leguas}.

%%K t09-c1-15-2 p
\VUgras{Mojones} \textsuperscript{(36)} e \VUgras{indicadores} \textsuperscript{(37)} jalonan
los caminos, en los que se encuentra a menudo a un joven
\VUgras{montañés} \textsuperscript{(38)} con la \VUgras{alforja} \textsuperscript{(40)} al
costado, que lleva una \VUgras{marmota} \textsuperscript{(39)}, o a algún
\VUgras{gitano} \textsuperscript{(41)} cuyo \VUgras{gorro de piel} \textsuperscript{(42)}
contrasta de un modo singular con su viejo \VUgras{capuchón}
\textsuperscript{(43)} roto. Lleva de la \VUgras{cadena} a un \VUgras{oso}
\textsuperscript{(44)} amaestrado.

%%K t09-tit-7 sub
\VUpk{10.2pt}{40}{Escalar montañas.}
\VUsaut{3.0mm}

%%K t09-c1-16-1 p
16. --- Casi todos los días voy con un \VUgras{guía} \textsuperscript{(48)} a hacer
una \VUgras{ascensión}, y estoy muy orgulloso cuando, con la \VUgras{mochila}
\textsuperscript{(47)} a la espalda y el \VUgras{bastón alpino} en la mano, como un
verdadero \VUgras{turista} \textsuperscript{(46)}, escalo las \VUgras{rocas}
\textsuperscript{(45)}.

%%K t09-c1-17-1 p
17. --- Desde lo alto de una montaña, los habitantes de la llanura no
parecen más grandes que hormigas; un \VUgras{rebaño de ovejas}
\textsuperscript{(50)} que pace en un \VUgras{pasto} \textsuperscript{(49)} parece un
juguete de niño. Un \VUgras{pino} \textsuperscript{(51)}, o incluso una \VUgras{cabaña
de madera} \textsuperscript{(52)}, apenas son una mancha en el verdor.

%%K t09-c1-18-1 p
18. --- Durante estas excursiones nos cruzamos a menudo en el
\VUgras{sendero} \textsuperscript{(53)} con un \VUgras{cazador de rebecos}
\textsuperscript{(54)} que lleva al hombro el animal que acaba de matar.

%%K t09-c1-19-1 p
19. --- He puesto en mi herbario una rama de \VUgras{helecho} \textsuperscript{(55)}
muy curiosa y un bonito ramito rosado de \VUgras{brezo} \textsuperscript{(57)}, que
cogí a la entrada de una pequeña \VUgras{gruta} \textsuperscript{(56)} en mi último
paseo.

%%K t09-tit-8 sub
\VUsaut{3.0mm}
\VUcentre{12.0pt}{0}{\textit{Segunda escena.}}
\VUsaut{3.0mm}

%%K t09-tit-9 sub
\VUpk{11.4pt}{40}{El Bosque. --- La Caza.}
\VUsaut{3.0mm}

%%K t09-c2-01-1 p
1. --- Ayer pasé un día delicioso en el bosque. La actividad que reina
entre los animales y los \VUgras{insectos} \textsuperscript{(5)} que lo pueblan me
interesó vivamente.

%%K t09-c2-01-2 p
La \VUgras{araña} \textsuperscript{(1)} que teje su \VUgras{tela} \textsuperscript{(2)} entre
dos ramas para atrapar a la \VUgras{mosca} \textsuperscript{(3)} que pasa, el
\VUgras{caracol} \textsuperscript{(4)} que lleva penosamente su concha, el ciervo
volante que busca su presa, la hormiga diligente que se afana junto al
\VUgras{hormiguero} \textsuperscript{(6)}: todo ese mundillo iba, venía y trabajaba
con una animación extraordinaria.

%%K t09-c2-02-1 p
2. --- Una \VUgras{serpiente} \textsuperscript{(7)}, que al principio tomé por una
\VUgras{víbora} pero que no era más que una simple \VUgras{culebra}, estaba
agazapada bajo un tronco de árbol y de vez en cuando sacaba su cabecita
fina, que siempre parecía dispuesta a morder. Delante de mí, una gruesa
\VUgras{babosa} \textsuperscript{(8)} reptaba pesadamente después de devorar una de
las sabrosas \VUgras{fresas silvestres} \textsuperscript{(11)} que crecen aquí en
abundancia.

%%K t09-c2-03-1 p
3. --- El suelo estaba alfombrado de \VUgras{flores del bosque}:
\VUgras{prímulas} \textsuperscript{(9)}, \VUgras{vincapervincas} \textsuperscript{(10)} y
muchas otras plantas que yo no conocía florecían entre las \VUgras{matas de
hierba} \textsuperscript{(12)}.

%%K t09-c2-04-1 p
4. --- En esta comarca, la \VUgras{trufa} es casi tan común como la
\VUgras{seta} \textsuperscript{(14)}, y un \VUgras{cerdito} \textsuperscript{(13)} de un
guarda forestal buscaba con aire goloso a ver si encontraba alguna.

%%K t09-tit-10 sub
\VUsaut{3.0mm}
\VUpk{10.2pt}{40}{La caza furtiva.}
\VUsaut{3.0mm}

%%K t09-c2-05-1 p
5. --- Asistí al \VUgras{final} de una escena que me divirtió mucho. Gracias
al olfato de su \VUgras{perro basset} \textsuperscript{(15)}, el \VUgras{guarda
forestal} \textsuperscript{(16)} acababa de sorprender a un \VUgras{cazador furtivo}
\textsuperscript{(22)} en el momento en que este se disponía a llevarse la
\VUgras{caza} \textsuperscript{(17)} que había matado. Prefiriendo abandonar su presa
a dejarse detener, el furtivo había huido, y \VUgras{liebre} \textsuperscript{(18)},
\VUgras{cierva} \textsuperscript{(19)}, \VUgras{corzo} \textsuperscript{(20)} y
\VUgras{jabalí} \textsuperscript{(21)} quedaban sobre la hierba, para alegría del
guarda.

%%K t09-c2-06-1 p
6. --- La variedad de árboles que encierra el bosque es asombrosa. Las
\VUgras{hayas} \textsuperscript{(23)} y los \VUgras{abedules} \textsuperscript{(26)}, los
\VUgras{pinos}, que dan la \VUgras{resina} \textsuperscript{(27)}, los \VUgras{robles}
\textsuperscript{(29)} y muchos más entremezclan sus ramas.

%%K t09-c2-06-2 p
La \VUgras{hiedra} \textsuperscript{(24)}, que se agarra al tronco de los árboles
altos, da al \VUgras{monte alto} \textsuperscript{(25)} un verde brillante que se une
agradablemente al tono más claro y más suave del \VUgras{musgo}
\textsuperscript{(31)} en el que el \VUgras{pito real} \textsuperscript{(30)} busca insectos.

%%K t09-tit-11 sub
\VUsaut{3.0mm}
\VUpk{10.2pt}{40}{Paisaje de bosque.}
\VUsaut{3.0mm}

%%K t09-c2-07-1 p
7. --- Los viejos robles me encantaron. Sus gruesas \VUgras{raíces}
\textsuperscript{(32)} retorcidas, medio salidas de la tierra, les dan una
espléndida apariencia de fuerza y de vigor, y me pregunto cómo el
\VUgras{propietario} \textsuperscript{(35)} puede resignarse a hacerlos talar y a
entregarlos a la \VUgras{sierra} \textsuperscript{(34)} del \VUgras{leñador}
\textsuperscript{(33)}. Los pobres \VUgras{troncos} \textsuperscript{(36)}, despojados de su
\VUgras{corteza} \textsuperscript{(37)} o hendidos por el \VUgras{hacha}
\textsuperscript{(38)} del \VUgras{jornalero} \textsuperscript{(39)}, me dan pena.

%%K t09-c2-08-1 p
8. --- Cerca de allí, un viejo \VUgras{carbonero} \textsuperscript{(41)} había
preparado con su \VUgras{pala} un \VUgras{montón} \textsuperscript{(42)} de carbón, y
una anciana se llevaba un \VUgras{haz} \textsuperscript{(40)} de \VUgras{leña muerta}
hecho con las ramillas de los árboles talados.

%%K t09-tit-12 sub
\VUsaut{3.0mm}
\VUpk{10.2pt}{40}{La Cacería.}
\VUsaut{3.0mm}

%%K t09-c2-09-1 p
9. --- Quería ir a descansar un rato en la \VUgras{casa del guarda forestal}
\textsuperscript{(46)}, pero cuando llegué junto al \VUgras{estanque} \textsuperscript{(43)},
en el que nada un hermoso \VUgras{cisne} \textsuperscript{(45)}, me di cuenta de que
una \VUgras{inundación} \textsuperscript{(44)} reciente había convertido el camino en
un verdadero \VUgras{pantano}. Tuve que dar un rodeo y, al pasar junto al
\VUgras{cedro} \textsuperscript{(49)}, los \VUgras{álamos} \textsuperscript{(48)} y el
\VUgras{olmo} \textsuperscript{(47)} que están en la linde del bosque, vi salir de un
\VUgras{soto} \textsuperscript{(54)} un magnífico \VUgras{ciervo} \textsuperscript{(50)}
perseguido por una \VUgras{jauría} \textsuperscript{(51)} encarnizada a la que
excitaba además el \VUgras{cuerno} \textsuperscript{(53)} de los \VUgras{monteros a
caballo} \textsuperscript{(52)}. El pobre animal estaba agotado. Vino a caer
moribundo a unos pasos de mí.

%%K t09-c2-10-1 p
10. --- El hijo del guarda forestal, atraído por el ruido de la
\VUgras{montería}, salió de la casa y volvimos los dos al bosque. Había
visto una \VUgras{urraca} \textsuperscript{(56)} en la rama de un viejo roble.
Pensaba que su nido debía de estar escondido en el \VUgras{follaje}
\textsuperscript{(58)} y quería cogerlo.

%%K t09-c2-10-2 p
Ágil como una \VUgras{ardilla} \textsuperscript{(61)}, trepó de \VUgras{rama}
\textsuperscript{(55)} en rama, pero no encontró nada y solo consiguió hacer volar
a una vieja \VUgras{corneja} \textsuperscript{(60)} posada en un \VUgras{ramo}
\textsuperscript{(57)}.
\VUsaut{4.0mm}
\VUfilet{22mm}
